%%% FIELD proof-adaptive-score-hull
Define
\[
 A_{n,P}=\left\{
 \sup_{\theta\in\Theta_I(P)}
 \frac{\sqrt n\{\widehat q_n(\theta)-q_P(\theta)\}}
 {\widehat\sigma_n(\theta)\vee n^{-1}}
 \leq\widehat c_{1-\eta}\right\}. \tag{1}
\]
For \(\theta\in\Theta_I(P)\), the sharp score characterization gives \(q_P(\theta)\leq0\). On \(A_{n,P}\),
\[
 \frac{\sqrt n\widehat q_n(\theta)}
 {\widehat\sigma_n(\theta)\vee n^{-1}}
 \leq
 \frac{\sqrt n\{\widehat q_n(\theta)-q_P(\theta)\}}
 {\widehat\sigma_n(\theta)\vee n^{-1}}
 \leq\widehat c_{1-\eta}. \tag{2}
\]
Thus \(\theta\in C_n\), and \(A_{n,P}\subseteq\{\Theta_I(P)\subseteq C_n\}\). Assertion (A) of the already-established uniform multiplier-coupling theorem now gives
\[
 \liminf_{n\to\infty}\inf_{P\in\mathcal P_W}
 P\{\Theta_I(P)\subseteq C_n\}\geq1-\eta. \tag{3}
\]

For loss, take \(\mathcal E_{n,P}\) and \(R_{n,P}\) from assertions (B)--(C) of that theorem. The endpoint-domain lemma shows that \(\Theta_I(P)\neq\varnothing\) for \(P\in\mathcal P_W\). Since both \(C_n\) and \(\Theta_I(P)\) are subsets of compact \(K\), the empty-set convention implies
\[
 0\leq d_{\to}\{C_n,\Theta_I(P)\}\leq\operatorname{diam}(K) \tag{4}
\]
on every sample. On \(\mathcal E_{n,P}\), assertion (C) gives
\[
 C_n\subseteq B_{n,P}:=\{\theta\in K:q_P(\theta)\leq R_{n,P}\}. \tag{5}
\]
Apply the score-error-set modulus pathwise with \(\widetilde q=q_P\) and \(r=R_{n,P}\); its approximation premise holds because the error is zero. Monotonicity of directed loss in its first argument under the inclusion (5) then gives
\[
 d_{\to}\{C_n,\Theta_I(P)\}
 \leq d_{\to}\{B_{n,P},\Theta_I(P)\}
 \leq\frac{\sqrt{\delta(P)+2R_{n,P}}-\sqrt{\delta(P)}}{c_\beta}. \tag{6}
\]
For \(x,u\geq0\),
\[
 \sqrt x\{\sqrt{x+2u}-\sqrt x\}\leq u,
 \qquad
 \sqrt{x+2u}-\sqrt x\leq\sqrt{2u}. \tag{7}
\]
The first inequality follows by rationalization and the second by squaring. Set \(a_n=n^{-1/4}\). Multiplying (6) by \(\sqrt{\delta(P)}+a_n\), applying (7), and then using (4) on \(\mathcal E_{n,P}^c\), we obtain
\[
 \begin{aligned}
 &\{\sqrt{\delta(P)}+a_n\}
 E_Pd_{\to}\{C_n,\Theta_I(P)\}\\
 &\quad\leq c_\beta^{-1}E_PR_{n,P}
 +\sqrt2c_\beta^{-1}a_nE_P\sqrt{R_{n,P}}\\
 &\qquad
 +\{\sqrt{\delta(P)}+a_n\}\operatorname{diam}(K)
 P(\mathcal E_{n,P}^c). \tag{8}
 \end{aligned}
\]
For \(P\in\mathcal P_W\), nonemptiness gives \(0\leq\delta(P)\leq\tau^2\), since \(\lambda(P)\geq0\). Uniformly in \(P\), the coupling theorem gives \(E_PR_{n,P}=O(n^{-1/2})\) and \(P(\mathcal E_{n,P}^c)=O(n^{-2})\); Jensen's inequality gives
\[
 E_P\sqrt{R_{n,P}}\leq\sqrt{E_PR_{n,P}}=O(n^{-1/4})
\]
Therefore the right-hand side of (8) is \(O(n^{-1/2})\), proving the continuous elbow bound.

If \(P\in\mathcal P_n^{\mathrm{ct}}(C_\delta)\), the multiplier \(\sqrt{\delta(P)}+n^{-1/4}\) is at least \(n^{-1/4}\), so the elbow bound gives \(O(n^{-1/4})\) loss. If \(P\in\mathcal P^{\mathrm{int}}(\delta_0)\), that multiplier is at least \(\sqrt{\delta_0}\), so it gives \(O(n^{-1/2})\) loss. The definition of \(C_n\), including its score statistic and multiplier critical value, is unchanged in these deductions, which proves adaptation by one rule.
